\documentclass[11pt,a4paper]{article}
\usepackage{fontspec}
\setmainfont{Noto Serif}
\setsansfont{Noto Sans}
\usepackage{babel}
\babelprovide[main,import]{russian}
\usepackage{amsmath,amssymb,mathtools,geometry,fancyhdr,microtype,xcolor,hyperref}
\geometry{left=21mm,right=21mm,top=20mm,bottom=20mm}
\definecolor{ink}{HTML}{173D57}
\hypersetup{colorlinks=true,urlcolor=ink,linkcolor=ink,pdftitle={BGS-F. Когерентные истории переходов}}
\pagestyle{fancy}\fancyhf{}\lhead{\small BGS-F}\rhead{\small Когерентные истории переходов}\cfoot{\thepage}
\setlength{\headheight}{15pt}\setlength{\parskip}{3pt}
\setlength{\emergencystretch}{3em}\tolerance=3000
\newcommand{\Tr}{\operatorname{Tr}}
\newcommand{\trn}{\operatorname{tr}_N}
\newcommand{\Ad}{\operatorname{Ad}}
\newcommand{\Op}{\operatorname{Op}}
\newcommand{\dist}{\operatorname{dist}}
\newcommand{\osc}{\operatorname{osc}}
\newcommand{\T}{\mathbb{T}}
\newcommand{\Z}{\mathbb Z}
\newcommand{\C}{\mathbb C}
\newcommand{\sect}[1]{\section*{\sffamily\color{ink}#1}}
\newcommand{\subt}[1]{\par\medskip\noindent\textbf{#1.}\quad}
\begin{document}
\raggedright
\noindent{\small\sffamily\color{ink} МАТЕМАТИЧЕСКАЯ ФИЗИКА\hfill 22 НОЯБРЯ 2026}
\begin{center}{\LARGE\sffamily Когерентные истории переходов}\par\vspace{4pt}
{\large Разделение квантовых амплитуд для точной карты BGS-F}
\end{center}

Пусть \(U_N=e^{i\tau NA_N}M_N\) при \(\tau=10^{-80}\). Берутся нечётные \(N,L\) с \(L^2\mid N\), \(L\ge3\). Обозначим через \(P_\chi\) спектральные проекторы группы \(L^2\) торсионных сдвигов Вейля, через \(\trn=\Tr/N\) нормированный след. Положим \(D=L^2\), \(d_L=16+\lceil\log L/\log 5\rceil\). Наблюдения выполняются в моменты \(0,d_L,2d_L,\ldots\).

\subt{Многократные истории}
С некоторой положительной постоянной \(c\) допустимы времена до \(T_N=\lfloor c\log\log\log N\rfloor\). Для \(q\) наблюдений положим \(H_*=D^q\), \(\gamma=(\chi_0,\ldots,\chi_{q-1})\), \(C_\gamma=\prod_{i=0}^{q-1} U_N^{i d_L}P_{\chi_i}U_N^{-i d_L}\), с умножением слева направо по возрастанию \(i\). Равномерно по историям и допустимым временам
\[
         \trn(C_\delta^* C_\gamma)
           =H_*^{-1}\mathbf 1_{\delta=\gamma}+O(\epsilon_N),
       \qquad \epsilon_N=\exp[-(\log N)^{3/4}].
\]
Оценка сохраняет фазы: она относится к скалярному произведению разных амплитуд, включая истории с общими началом и концом. В частности, средние условные вероятности последовательных наблюдений асимптотически равны \(D^{-(q-1)}\).

\subt{Следы любых степеней}
Если в допустимое окно помещаются \(s\) непересекающихся копий блока из \(q\) наблюдений, то для всех \(m\in\Z\)
\[
 H_*\sum_\gamma\left|\trn(C_\gamma U_N^m)
             - H_*^{-1}\trn(U_N^m)\right|^2
       \le (1+O(sH_*^2\epsilon_N))/s .
\]
Рост \(s\) обеспечивает распределение нормированного следа по историям в среднем квадрате и на временах \(m\asymp N\).

\subt{Истории в собственных состояниях}
При фиксированном \(L\) можно одновременно взять \(q\to\infty\) и \(s\to\infty\) так, что в любой ортонормированной системе собственных состояний полные законы наблюдений у доли \(1-o(1)\) состояний сходятся по полной вариации к равномерному закону на \(D^q\) историях.

\subt{Статус задачи}
Явный пример с другим унитарным оператором удовлетворяет тем же тождествам для историй точно, а его парный тест расходится. Усиленный пример сохраняет контроль историй, простоту и независимость фаз и подавляет все субгейзенберговские следы, но даёт неверный микроскопический закон. У исходной карты общий скалярный след, парный тест \((P)\) и предел \(\mathrm{Sine}_2\) остаются открытыми.

\newpage
\sect{1. Точные группы и произведения проекторов}
Для \(N\) нечётного обозначим \(e_N(u)=e^{2\pi i u/N}\). В исходных координатах
\[
 (T_N(a,b)v)(x)=e_N(b(x+\overline2 a))\,v(x+a),\qquad
 (M_N)_{xy}=N^{-1/2}e_N(-\overline8(3x^2-2xy+3y^2)).
\]
Кик \(A_N\) и его физический символ \(g\) заданы в предыдущей главе. Ни величина кика, ни метаплектическая фаза здесь не изменяются. Точная ковариантность даёт \(M_N T_N(a,b)M_N^*=T_N(S(a,b))\), где \(S=\bigl(\begin{smallmatrix}3&4\\ 2&3\end{smallmatrix}\bigr)\).

Пусть \(\Gamma_L=(L^{-1}\Z/\Z)^2\). При \(a=(-u/L,v/L)\) используем обозначение \(W_a=T_N(Nu/L,Nv/L)\). Это точное представление \(\Gamma_L\), поскольку в центральной фазе произведения появляется множитель \((N/L)^2\), делящийся на \(N\). След \(\trn W_a\) равен \(\mathbf 1_{a=0}\). Формулы для проекторов и их переносов таковы:
\[
 P_\chi=D^{-1}\sum_{a\in\Gamma_L}\overline{\chi(a)}W_a,\qquad
 W^t_a=U_N^t W_a U_N^{-t},\qquad
 P^t_\chi=U_N^t P_\chi U_N^{-t}.
\]
Характеры \(\chi\) являются групповыми гомоморфизмами \(\Gamma_L\to\mathbb S^1\). Матрицы \(W^t_a\) при фиксированном \(t\) умножаются сложением меток. Проекторы \(P_\chi^t\) имеют ранг \(N/D\), ортогональны попарно и в сумме дают \(I\).

\subt{Последовательные наблюдения}
Для сетки \(t_j=t_0+j d_L\), \(0\le j<q\), положим
\[
 C_\gamma(t_0)=P^{t_0}_{\chi_0}\cdots P^{t_{q-1}}_{\chi_{q-1}}.
\]
При \(t_0=0\) матрица \(C_\gamma(0)U_N^{t_{q-1}}\) есть произведение
\(P_{\chi_0}U_N^{d_L}P_{\chi_1}\cdots U_N^{d_L}P_{\chi_{q-1}}\).
Так оно действует на вход из сектора \(\chi_{q-1}\), а наблюдения идут от правого фактора к левому. Средняя вероятность заданной истории при равномерной смеси во входном секторе равна \(D\,\trn(C_\gamma^* C_\gamma)\). Для собственного вектора \(v\) исторические вероятности равны \(\|C_\gamma v\|^2\). Для любой сетки в точности
\[
               \sum_\gamma C_\gamma=I,\qquad
               \sum_\gamma C_\gamma^* C_\gamma=I .
\]
Первое тождество получается прямым суммированием, второе последовательным удалением соседних проекторов одинакового времени после суммирования.

\subt{План доказательства}
Классические отображения \(F^t t_a F^{-t}\) задают переносы \(W_a^t\) в формуле Егорова. Изолируем слова с тождественной классической картой в трёх нужных типах произведений. Для остальных получим оценку \(|\trn(\text{слово})|\le\epsilon_N\), с константами, допускающими растущее число факторов. Разложения по характерам затем дают точные главные члены.

\newpage
\sect{2. Наибольшее время в произведении сдвигов}
Классическая карта \(F=\varphi^\tau B\), где \(B=\bigl(\begin{smallmatrix}3&-4\\-2&3\end{smallmatrix}\bigr)\), \(\varphi^t\) — поток поля \(Y=(-g_p,g_q)/(2\pi)\). Положим \(\lambda=3+2\sqrt2\), \(\ell=(1/2,-1/\sqrt2)\), тогда \(\ell B=\lambda\ell\). Пусть \(F H=HB\), \(H^{-1}=\mathrm{id}+k\). Используем точные оценки \(\|H-\mathrm{id}\|_\infty,\|k\|_\infty\le100\tau\) и
\[
           |\widehat{\ell k}(B^{\mathsf T}e_i)|>10^{-4}\tau
                \qquad(i=1,2).
\]
Они уже использованы в главах о каналах Вейля и сдвинутых возвратах. В разделе 9 приведена проверка коэффициентов и сопряжения.

Для \(t_a(x)=x+a\) сопряжённый сдвиг \(H^{-1}t_a H\) имеет поднятие со смещением \(a+k(Hx+a)-k(Hx)\). Обозначим через \(\osc\) разность супремума и инфимума. Так как \(H\) сюръективно, осцилляция \(\ell\)-компоненты совпадает с осцилляцией \(\ell(k(\cdot+a)-k)\). Для \(a\ne0\) одна из двух указанных мод имеет амплитуду не меньше \(4\cdot10^{-4}\tau\,\|a\|_{\T,\infty}/7\): частоты образуют унимодулярный базис с матрицей \(B\), а \(\|B^{-1}\|_\infty=7\). Модуль нетривиального коэффициента не превосходит осцилляции. Верхняя граница осцилляции получается из \(\|k\|_\infty\). Поэтому для
\(G^t_a=B^t(H^{-1}t_a H)B^{-t}\), \(a\in\Gamma_L\setminus\{0\}\), имеем
\[
    10^{-5}\tau\lambda^t/L
       \le \osc\bigl(\ell(G^t_a-\mathrm{id})\bigr)
       \le 1000\tau\lambda^t .                               \tag{1}
\]
Осцилляция не зависит от выбора поднятия.

Смещение композиции равно сумме смещений факторов, вычисленных на хвостовых композициях. Хвостовые карты сюръективны, поэтому осцилляции каждого слагаемого сохраняются. Пусть в слове старшее ненулевое время \(t\) встречается один раз, а каждое из меньших времён не более двух раз и расстояния между временами не меньше \(d_L\). Тогда
\[
 \osc\bigl(\ell(\text{композиция}-\mathrm{id})\bigr)
 \ge \frac{10^{-5}\tau\lambda^t}{L}
             -\frac{2000\tau\lambda^t}{\lambda^{d_L}-1}
 \ge \frac{10^{-5}\tau\lambda^t}{2L}.                          \tag{2}
\]
В частности, при \(t\ge d_L\) нижняя граница больше \(10^5\tau\), поскольку \(\lambda>5\) и \(5^{16}>10^{11}\). Если исходное слово в картах \(F^t t_a F^{-t}\) было равномерно ближе \(\tau\) к тождеству на торе, то сопряжённое с помощью \(H\) слово было бы ближе \(201\tau\). У его поднятия \(\ell\)-осцилляция тогда меньше \(600\tau\): ближайший целый вектор к смещению постоянен по связности тора. Это противоречит (2).

\subt{Удаление центральных факторов}
В слове с временами \(t_0,t_1,\ldots,t_r,\ldots,t_1,t_0\) соседние факторы одного старшего времени сначала умножаются, а результат с нулевой меткой удаляется. Процедура повторяется. Если слово дошло до ненулевого старшего фактора, применима (2). Для произведения двух таких слов на непересекающихся временных блоках сначала удаляем центральные факторы блока с большими временами. Как только остался ненулевой старший фактор, оценки (1), (2) охватывают все оставшиеся меньшие времена, включая другой блок. Если старшее время равно нулю, остаётся обычный сдвиг, нормированный квантовый след которого точно равен нулю.

\newpage
\sect{3. Количественная оценка следа целого слова}
В нужных словах суммарное число шагов \(U_N^{\pm1}\) после объединения смежных степеней не превосходит \(4T\), где все времена лежат в \([0,T]\). Достаточно сложить абсолютные разности соседних времён, включая начальный и конечный переход к нулю. У отдельных слов времена идут вверх и вниз, а у произведения двух слов на непересекающихся блоках имеется не более двух подъёмов. Факторы с одинаковым временем умножаются без фазы.

\subt{Аналитическая область}
Локально в комплексной трубке фиксированной ширины \(r_0>0\) подъёмы \(F,F^{-1}\) аналитичны, вещественны на вещественной плоскости и увеличивают норму мнимой части не более чем в \(K>1\) раз. Это следует из конечного тригонометрического символа кика, локального аналитического потока и интегрирования производной по вертикальному отрезку. Вещественные сдвиги трубку сохраняют. Поэтому классическое слово \(G\) аналитично на трубке \(r_0 K^{-4T}\). Его мнимая часть там ограничена \(r_0\). По (2) для нетривиального слова есть координата и точка \(x_0\) с
\[
           |f(x_0)|\ge 4\tau,\qquad
        f(z)=e^{2\pi i(G_j(z)-z_j)}-1.
\]
Модуль \(f\) на трубке ограничен фиксированной постоянной \(M\).

\subt{Мера области малых смещений}
Сократим ширину трубки до \(\rho=r_0K^{-4T}/2\). На каждом вещественном луче из \(x_0\), до длины 2, голоморфно переведём полосу \(|\Im z|<\rho\) в диск через \(\zeta=\tanh(\pi z/(4\rho))\). Образ отрезка заканчивается в \(Q=\tanh(\pi/(2\rho))\), положим \(b=1-Q\). По формуле Йенсена число нулей на диске радиуса \(Q+b/3\) не превосходит \(C_\tau/b\), сравнивая с радиусом \(Q+2b/3\). Удалим эти нули множителями Бляшке. Безнулевой остаток ограничен сверху \(M\); по Харнаку на \([0,Q]\) его модуль не меньше \(M e^{-C_\tau/b}\). Радиусы можно выбирать на расстоянии сравнимом с \(b\) без нулей на окружностях.

Если \(\varepsilon/M<e^{-2C_\tau/b}\), то при \(|f|\le\varepsilon\) произведение множителей не больше \((\varepsilon/M)^{1/2}\). Поэтому точка \(\zeta\) находится на расстоянии \(O((\varepsilon/M)^{c_\tau b})\) от одного из нулей. Обратное отображение на вещественном отрезке имеет производную \(O(1/b)\). Длина множества не больше \(C_\tau b^{-2}\varepsilon^{c_\tau b}\). После увеличения постоянной оценка верна для остальных \(\varepsilon\le1\). Интегрирование по лучам до длины 2 покрывает фундаментальную ячейку с ограниченным радиальным якобианом. Так как \(b^{-1}\le C e^{C K^{4T}}\), получаем с подходящими постоянными
\[
 \big|\{x:\dist(G(x),x)\le\varepsilon\}\big|
   \le C_\tau e^{C_\tau K^{4T}}\,
                    \varepsilon^{e^{-C_\tau K^{4T}}}.             \tag{3}
\]

\subt{Квантовая оценка}
Подробности перехода вынесены на следующую страницу. В результате, если
\(T\le T_N=\lfloor c\log\log\log N\rfloor\) при достаточно малом \(c>0\),
\[
     |\trn(\text{нетривиальное слово})|\le \epsilon_N .          \tag{4}
\]
Если процедура удаления приводит к пустому слову, в квантовой задаче получается точно \(I\).

\newpage
\sect{4. Перенос оценки с тора на матрицы}
В обычной конвенции \(B_\nu=(-1)^{kl}T_N(l,k)\) при \(\nu=(k,l)\), а
\(\Op_N(f)=\sum_\nu\widehat f(\nu)B_\nu\). Точная фаза в произведении \(B_\nu B_\mu\) равна \(e^{\pi i(l k'-k l')/N}\) относительно \(B_{\nu+\mu}\). Поэтому при гладких \(a,b\)
\[
  \|\Op_N(a)\Op_N(b)-\Op_N(ab)\|
        \le C N^{-1}\|a\|_{C^6}\|b\|_{C^6}.
\]
Если один символ есть одиночная гармоника с частотой \(\le R\), вместо его \(C^6\)-нормы достаточно \(O(1+R)\) по той же фазовой формуле.

У генератора \(i\tau N A_N\) коммутатор даёт классическую производную с ошибкой \(O(N^{-2}\|a\|_{C^8})\) после линеаризации \(\sin x\). Это следует из конечного набора частот \(g\) и оценки взвешенной суммы \(\sum(1+|\nu|)^3|\widehat a(\nu)|\le C\|a\|_{C^8}\). Формула Дюамеля даёт тот же порядок для одного кика или обратного кика. Сопряжение линейной частью и квантованными сдвигами точно. Нормы переноса \(a\mapsto a\circ F^{\pm1}\) в \(C^8\) ограничены фиксированной постоянной, а сдвиги — изометрии. Следовательно, для кванта \(V\) слова \(G\)
\[
       V^*\Op_N(a)V
           =\Op_N(a\circ G)+O(C^{T+1}N^{-2}\|a\|_{C^8}).
\]

Усредним \(B_\nu^* V B_\nu\) по \(|\nu|_\infty\le R\) в матрицу \(J_R\). Все следы равны. После домножения слева на \(V^*\) получаем
\[
 V^*J_R=\Op_N(a_R)+O(C^{T+1}N^{-1}(1+R)^8),\qquad
 a_R(x)=D_R(x_1-G_1(x))D_R(x_2-G_2(x)),
\]
где \(D_R(u)=(2R+1)^{-1}\sum_{j=-R}^R e^{2\pi iju}\).
В норме \(\|A\|_{2,N}^2=\trn(A^*A)\) справедливо
\(\|\Op_N(a)\|_{2,N}\le\|a\|_{L^2}+CN^{-1/2}\|a\|_{C^6}\).
Это следует из оценки произведения и следового правила: кроме нулевой частоты в следе остаются только частоты, кратные \(N\) по обеим координатам, с суммарной ошибкой \(O(N^{-4}\|\,|a|^2\|_{C^4})\). С \(\|a_R\|_{C^6}\le C^{T+1}(1+R)^6\) выводим
\[
   |\trn V|\le \|a_R\|_2+
       C^{T+1}\bigl(N^{-1/2}(1+R)^6+N^{-1}(1+R)^8\bigr).
\]
Вне множества из (3) при \(\varepsilon=R^{-1/2}\) одна из геометрических сумм имеет модуль \(O(R^{-1/2})\), другая не больше 1. Берём \(R=\lfloor N^{1/32}\rfloor\). Если \(c\) настолько мало, что \(K^{4T_N}\le(\log\log N)^{1/4}\), оценка (3) с этой \(\varepsilon\) меньше \(\exp[-2(\log N)^{3/4}]\) с запасом в показателе. Операторные остатки убывают степенно по \(N\). Это доказывает (4) равномерно по словам, \(L\) и характерам.

\newpage
\sect{5. Ортогональность историй}
Пусть \(q,s\ge1\), \((sq-1)d_L\le T_N\) и \(d_L\le T_N\).
Расположим \(s\) блоков из \(q\) времён на общей сетке, начиная с нуля. Пишем \(C_\gamma^{[b]}=C_\gamma(bqd_L)\), \(0\le b<s\), \(H_*=D^q\). Всюду \(\epsilon_N=\exp[-(\log N)^{3/4}]\). Докажем три формулы:
\begin{align}
 \trn C_\gamma^{[b]}&=H_*^{-1}+O(\epsilon_N), \tag{5}\\
 \trn\!\big((C_\delta^{[b]})^*C_\gamma^{[b]}\big)
     &=H_*^{-1}\mathbf1_{\delta=\gamma}+O(\epsilon_N),\tag{6}\\
 \trn\!\big((C_\delta^{[b']})^*C_\gamma^{[b]}\big)
     &=H_*^{-2}+O(\epsilon_N)\qquad(b\ne b').\tag{7}
\end{align}
Ошибка во всех развёртках контролируется (4): сумма модулей весов в полном разложении проекторов на сдвиги равна 1.

В (5) каждое время встречается ровно один раз. Если есть ненулевая метка, её старшее время выделяется по (2) или даёт чистый сдвиг. В главном члене все метки равны нулю. В (7) все времена также различны, поэтому главный член получается таким же способом. Для (6) используем цикличность следа и записываем времена по возрастанию до общей вершины и обратно. Обозначим метки двух историй \(a_i,a_i'\). Сложение в вершине и удаление нулевых факторов последовательно показывают, что тождественный вклад возможен в точности при \(a_i'=-a_i\) для всех \(i\). Сумма фаз на каждом времени даёт ортогональность характеров, и коэффициент равен \(H_*^{-1}\mathbf1_{\delta=\gamma}\). Это обосновывает (5)--(7), включая исходные случаи с одинаковыми начальными или конечными секторами.

\subt{Векторный вывод}
Для одного блока и любых комплексных \(c_\gamma\) из (6) получаем
\[
  \left\|\sum_\gamma c_\gamma C_\gamma\right\|_{2,N}^2
     =H_*^{-1}\bigl(1+O(H_*^2\epsilon_N)\bigr)
                               \sum_\gamma|c_\gamma|^2 .
\]
Можно брать \(q\) и \(L\) переменными. По выбору \(d_L\) и \(d_L\le T_N\) имеем
\(\log H_*=2q\log L\le4(\log 5) T_N\), так что \(sH_*^2\epsilon_N\to0\).
В частности, разложение по всем историям одновременно почти ортогонально. Это утверждение о комплексных матричных амплитудах.

\subt{Отрицательный контроль при нулевом кике}
Для \(U_N=M_N\) метаплектическая ковариантность переставляет проекторы по характерам. Исходный сектор определяет целую историю. Поэтому (6) неверна при \(\tau=0\) для \(q\ge2\): ненулевые истории имеют средний вес \(1/D\) вместо \(1/D^q\). Гиперболическое расширение в (2) использует ненулевые коэффициенты сопряжения.

\subt{Истории операторов Вейля по всем нечётным \(N\)}
Для этого следствия условие \(L^2\mid N\) не требуется. Возьмём множество \(E\) физических сеточных сдвигов, содержащее 0, с попарными расстояниями на торе не меньше \(1/L\) в максимальной норме. При тех же \(d_L,q,s\) обозначим через \(X_\gamma^{[b]}\) произведение перенесённых вейлевых операторов с метками \(\gamma\in E^q\) в порядке возрастания времени. Совпадающие метки в паре с сопряжённой историей сокращаются точно. Различные метки после умножения дают сдвиг на расстоянии \(\ge1/L\) от 0 и скалярную фазу модуля один. Для такого сдвига работают (1)--(4). Поэтому в одном блоке
\(\trn((X_\delta^{[b]})^*X_\gamma^{[b]})=\mathbf1_{\delta=\gamma}+O(\epsilon_N)\).
Между блоками главный член есть лишь у двух полностью нулевых наборов.

\newpage
\sect{6. Следы произвольных степеней и собственные состояния}
\subt{Равномерность для всей коммутирующей алгебры}
Возьмём \(Y_\gamma^{[b]}=\sqrt{H_*}(C_\gamma^{[b]}-H_*^{-1}I)\).
Из (5)--(7) матрица Грама этих \(sH_*\) векторов отличается по операторной норме на \(O(sH_*^2\epsilon_N)\) от блочно-диагональной матрицы с проектором \(I-{\bf1}{\bf1}^*/H_*\) в каждом блоке. Получаем неравенство Бесселя с правой константой \(1+O(sH_*^2\epsilon_N)\). Если \(Z\) коммутирует с \(U_N\), скалярные произведения \(\trn((Y_\gamma^{[b]})^*Z)\) не зависят от \(b\), поскольку копии связаны сопряжением \(U_N^{b q d_L}\). Следовательно,
\[
 H_*\sum_\gamma
    \left|\trn(C_\gamma^*Z)-H_*^{-1}\trn Z\right|^2
     \le \frac{1+O(sH_*^2\epsilon_N)}{s}\,\|Z\|_{2,N}^2 .        \tag{8}
\]
При \(Z=U_N^{-m}\) и комплексном сопряжении это даёт оценку со страницы 1. Она охватывает любое целое \(m\), поскольку используется точная коммутативность.

Для вейлевых историй из раздела 5 исключаем полностью нулевой набор меток. Их объединённая матрица Грама отличается от единицы на \(O(s|E|^q\epsilon_N)\). Тот же аргумент для \(Z U_N=U_N Z\) даёт на всех нечётных \(N\)
\[
      \sum_{\gamma\in E^q,\;\gamma\ne0}
          |\trn(X_\gamma^* Z)|^2
          \le \frac{1+O(s|E|^q\epsilon_N)}{s}\|Z\|_{2,N}^2 .
\]
Здесь \(|E|\le C L^2\) по упаковке. Это ограничивает коллективный вклад произведений вейлевых сдвигов в любой степени \(U_N\), в том числе при простых \(N\).

\subt{Законы наблюдений в собственных состояниях}
Обозначим \(\Omega_\gamma^{[b]}=(C_\gamma^{[b]})^*C_\gamma^{[b]}\). Из (6) \(\trn\Omega_\gamma^{[b]}=H_*^{-1}+O(\epsilon_N)\), а \(0\le\Omega_\gamma^{[b]}\le I\). При \(b\ne b'\) развёртка \(\trn(\Omega_\gamma^{[b]}\Omega_\gamma^{[b']})\) содержит два слова с непересекающимися наборами времён. В каждом слово идёт вниз к своему минимальному времени, а затем вверх.

Берём последнее время сетки \(T=(sq-1)d_L\) и сопрягаем всё произведение \(U_N^{-T}\). Теперь времена \(T-t\) относятся к обратной карте \(F^{-1}\). Блоки становятся словами с максимумом в середине. Из оценки \(|\widehat{\ell_s k}(e_i)|>10^{-4}\tau\) для устойчивого ковектора \(\ell_s=(1/2,1/\sqrt2)\), проверенной в разделе 9, вывод (1)--(4) работает тем же способом с \(\ell_s B^{-1}=\lambda\ell_s\). В каждой вершине остаётся сумма меток парных факторов на каждом времени. Старший блок редуцируется первым. Как и в (6), каждый блок даёт для тождества вес \(H_*^{-1}\). Итак,
\(\trn(\Omega_\gamma^{[b]}\Omega_\gamma^{[b']})=H_*^{-2}+O(\epsilon_N)\).

Для любой ортонормированной системы собственных состояний \((v_j)_{j=1}^N\) диагональные элементы \(\Omega_\gamma^{[b]}\) не зависят от \(b\). Усредняем \(\Omega_\gamma^{[b]}-H_*^{-1}I\) по блокам, берём квадрат нормы \(\|\cdot\|_{2,N}\) и сохраняем только диагональные элементы в собственном базисе. Для одного \(\gamma\) получаем средний квадрат \(\le1/(sH_*)+O(\epsilon_N)\). По Коши и суммированию
\[
 \frac1N\sum_j\sum_\gamma
       \big|\|C_\gamma v_j\|^2-H_*^{-1}\big|
                  \le \sqrt{H_*/s+O(H_*^2\epsilon_N)} .          \tag{9}
\]
Для фиксированного \(L\) и \(J=1+\lfloor T_N/d_L\rfloor\) можно взять
\(q=\lfloor(\log J)/(4\log D)\rfloor\) и \(s=\lfloor J/q\rfloor\) при больших \(N\). Тогда \(H_*\le J^{1/4}\) и правая часть стремится к нулю.

\newpage
\sect{7. Что остаётся в парном тесте}
Пусть \(p_{m,N}=\Tr U_N^m\). Точное тождество, следующее из \(\sum C_\gamma=I\), вместе с (8) даёт
\begin{align*}
 H_*\sum_\gamma|\trn(C_\gamma U_N^m)|^2
    &=|\trn(U_N^m)|^2+
       H_*\sum_\gamma\left|\trn(C_\gamma U_N^m)
                     -H_*^{-1}\trn U_N^m\right|^2,\\
 0&\le H_*\sum_\gamma|\trn(C_\gamma U_N^m)|^2
                      -|\trn(U_N^m)|^2
       \le \frac{1+O(sH_*^2\epsilon_N)}{s}.
\end{align*}
Если на некотором длинном времени нормированный след существенно отличен от нуля, он распределяется по типичным коротким историям. Однако верхняя ошибка после суммирования по гейзенберговской полосе с исходным весом имеет порядок \(N/s\). Даже при её усилении требовалось бы отдельно вычислить общий скалярный вклад.

\subt{Модель с точной статистикой историй и расходящимся тестом}
Зафиксируем \(L\), \(D=L^2\), возьмём \(N=D^B\) и представим
\(\C^N=(\C^D)^{\otimes B}\). Равномерные проекторы ранга \(N/D\) можно точным унитарным преобразованием отождествить с координатными проекторами первого множителя. Возьмём вспомогательный оператор \(Z_B\), циклически переставляющий \(B\) множителей. Проекторы в разные моменты \(0,\ldots,B-1\) тогда действуют на разные множители. Поэтому (5)--(9) выполняются для этой модели при \(T_N<B\) с нулевыми ошибками в следовых тождествах. Формулы верны даже на последовательной сетке с шагом 1.

Но \(Z_B^B=I\). На любом фиксированном отрезке \([a,b]\subset(0,2)\), где исходный вес \(w\) ограничен снизу положительной постоянной, множество \(m\) с \(m/N\in[a,b]\) и \(B\mid m\) имеет мощность порядка \(N/B\). При таких \(m\), \(N^{-2}|\Tr Z_B^m|^2=1\). Вычитаемая CUE-сумма ограничена, откуда парный дефект для \(Z_B\) стремится к \(+\infty\). Это уточняет предел применимости метода: даже точный контроль всех историй в наблюдаемом окне не определяет гейзенберговский форм-фактор.

\subt{Новый вклад и зависимости}
Ранее были установлены одношаговая статистика секторов, отделённость одиночного сдвига при многих итерациях и ограничения на отдельные торсионные следы. Здесь осцилляция старшего сдвига разделяет многократные слова, что даёт следовые тождества для когерентных историй, оценку для всей коммутирующей алгебры и результат для законов измерений в собственных состояниях. Применение к исходной карте при \(m\asymp N\) выражено именно в (8). Нужный предел полного парного теста \((P)\) остаётся открытым.

\subt{Проверка}
Скрипт \texttt{histories/check.py} проверяет численные константы, тождества сопряжения, удаление парных меток, конечные матричные формулы и тензорный контрпример. Аналитический предел доказывается на страницах 3--7. Конечные тесты служат для контроля формул, знаков и нормировок.

\newpage
\sect{8. Пример с контролем коротких следов и простыми фазами}
Усиление предыдущего примера. Берём \(L=3\), \(D=9\), \(B=3^j\), \(N=D^B\).
Циклическая перестановка \(B\) тензорных множителей переставляет слова длины \(B\). Длины её циклов \(h\) делят \(B\). Выберем фиксированное \(b_0=9^a\), \(a\ge1\), и положим \(K_0=N/b_0\). При больших \(j\) все \(h\) делят \(K_0\), и семейство подгрупп порядка \(h\) в \(\Z/K_0\) вложено по размерам.

Распределим по циклам перестановки \(b_0\) копий равномерной сетки \(\Z/K_0\). Циклу длины \(h\) приписывается неиспользованный косет подгруппы размера \(h\) внутри одной копии. Это можно делать по убыванию длин циклов. Занятые ранее большие косеты — объединения косетов текущего размера. Число оставшихся точек равно сумме оставшихся длин циклов и в частности не меньше \(h\). Значит, подходящий свободный косет существует на каждом шаге.

Теперь домножим оператор перестановки слева на диагональную унитарную матрицу. Для цикла длины \(h\), которому приписан косет с представителем \(u\), задаём произведение фаз по циклу равным \(e^{2\pi i u h/K_0}\). На таком цикле характеристический многочлен \(z^h-e^{2\pi i u h/K_0}\), поэтому его фазы в точности составляют приписанный косет. Для результата \(V_N\) имеем
\[
    \Tr V_N^m=N\mathbf1_{K_0\mid m}.
\]
Сопряжение диагональных координатных проекторов не изменилось, поэтому следовые тождества историй остаются точными на всех сетках между 0 и \(B-1\).

Выбираем \(b_0\) достаточно большим, чтобы сумма \(w(j'/b_0)\) по узлам одного фиксированного отрезка \([a_1,a_2]\subset(0,2)\) с \(\inf w>0\) превысила \(2\|w\|_\infty+1\). Это возможно по росту числа узлов. В парной сумме уже одни степени \(m=j'K_0\) с \(m/N\) на этом отрезке дают превосходство над равномерно ограниченной CUE-частью. При этом все ненулевые следы со степенью \(m<K_0\) исчезают. Здесь и далее рассматриваем положительные \(m\).

\subt{Простота и независимость}
Пусть \(p_1,\ldots,p_N\) — разные простые числа. В собственном базисе \(V_N\) добавим к фазам углы
\[
             \delta_i=\frac{\sqrt{p_i}-N^{-1}\sum_j\sqrt{p_j}}{2N^3 p_N}.
\]
Углы различны и по модулю не больше \(N^{-3}\). Полученная матрица \(V_N'\) сохраняет определитель и имеет простой спектр. Для целых \(n_i\) равенство \(\sum_i n_i\delta_i=0\) требует постоянства всех \(n_i\) по линейной независимости квадратных корней разных простых. По теореме Линдемана--Вейерштрасса произведение фаз с показателями \(n_i\) алгебраично только при постоянном наборе \(n_i\).

По телескопированию \(\|(V_N')^t-V_N^t\|\le |t|N^{-3}\).
Истории длины \(q\) на окне до \(T_N\) меняются в операторной норме на \(O(qT_N N^{-3})\), что меньше \(\epsilon_N\) с запасом. В парной сумме на временах \(m\le2N\) замена меняет нормированный след на \(O(N^{-2})\), а всю взвешенную квадратичную сумму на \(O(N^{-1})\). В частности, при \(1\le m<K_0\) ненормированные следы \(V_N'\) равномерно \(O(N^{-1})\). После обычной развёртки с равномерным сдвигом исходная сетка имеет шаг \(b_0\) и кратность \(b_0\). Добавки к развёрнутым фазам равномерно \(O(N^{-2})\), поэтому предел \(V_N'\) есть случайно сдвинутая решётка с той же кратностью, отличная от простого процесса \(\mathrm{Sine}_2\).

\newpage
\sect{9. Проверка двух направлений сопряжения}
Физический символ \(g=\cos(14\pi q)-1/10+r(q)(1/10+\sin(2\pi q)/100)
+r(q+1/2)\cos(2\pi p-\alpha_0)/100\), где
\(r(q)=(1-\cos(14\pi q))(1+\cos(2\pi q))\), \(e^{i\alpha_0}=(-3+4i)/5\).
Из его 32 коэффициентов получаем \(\sum|\widehat g|<1.2\),
\(\|Y\|_\infty<20,\|DY\|_\infty<10^3\).
Для смещения \(d=\varphi^\tau-\mathrm{id}\) интегральное уравнение потока даёт
\(\|d\|_\infty\le20\tau\), \(\mathrm{Lip}(d)\le2\cdot10^3\tau\),
\(\|d-\tau Y\|_\infty\le10^4\tau^2\).

Стабильный и нестабильный проекторы \(P_s,P_u\) матрицы \(B\) имеют нормы \(<2\). Уравнение \(h(x)-B h(B^{-1}x)=f(x)\) имеет обратную \(\mathcal L^{-1}\) через суммы
\[
 (\mathcal L^{-1}f)_s=\sum_{j\ge0}B^jP_s f(B^{-j}x),\qquad
 (\mathcal L^{-1}f)_u=-\sum_{j\ge1}B^{-j}P_u f(B^j x).
\]
Норма в \(C^0\) меньше 5. При \(f(x)=d(x+B h(B^{-1}x))\) это сжатие, поскольку \(5\cdot7\cdot\mathrm{Lip}(d)<1\). Значит, \(H=\mathrm{id}+h\) с \(\|h\|_\infty\le100\tau\) удовлетворяет \(FH=HB\). Сюръективность следует из степени 1. Если \(H x=H y\), все \(B^j x,B^j y\), \(j\in\Z\), остаются близки с точностью \(200\tau\). Центрированные разности при таких расстояниях подчиняются точной линейной рекурсии, потому что \(\|B\|_\infty=7\) исключает целочисленный скачок. Ограниченный во всех целых направлениях вектор гиперболической рекурсии нулевой. Получаем гомеоморфизм и \(\|k\|_\infty\le100\tau\).

Пусть \(h_1=\mathcal L^{-1}Y\). Из уравнения выше
\(\|h-\tau h_1\|_\infty\le10^8\tau^2\). Из двух сумм видно, что
\([h_1]_{1/2}<2\cdot10^4\): для ряда каждого направления используем
\(\|B^{\pm j}\|_\infty\le 2\lambda^j\) и суммируем \(\lambda^{-j/2}\).
Поскольку \(k=-h\circ H^{-1}\), получаем
\(\|k+\tau h_1\|_\infty\le10^8\tau^2+10^6\tau^{3/2}\).

Носитель \(\widehat g\) удовлетворяет \(|l|\le1\) во второй координате частоты \((k,l)\). Ни одна ненулевая степень \(B^{\mathsf T}\), применённая к \(e_1\) или \(e_2\), не попадает в этот носитель. Вторые координаты по модулю уже на шагах \(\pm1\) равны 4 и 3 соответственно и далее растут по рекурсии \(x_{j+1}=6x_j-x_{j-1}\).
В нестабильном ряде при частоте \(B^{\mathsf T}e_i\) остаётся ровно член \(j=1\), равный \(-\lambda^{-1}\ell\,\widehat Y(e_i)\). В стабильном ряде при частоте \(e_i\) остаётся член \(j=0\), равный \(\ell_s \widehat Y(e_i)\). Модули \(|\ell\widehat Y(e_i)|\), \(|\ell_s\widehat Y(e_i)|\) равны \(|\widehat g(e_1)|/\sqrt2\) для \(i=1\) и \(|\widehat g(e_2)|/2\) для \(i=2\). Здесь \(\widehat g(e_1)=1/20-i/200\), \(\widehat g(e_2)=e^{-i\alpha_0}/200\). Так как \(\lambda<6\), после учёта остатка обе пары коэффициентов \(k\) по модулю больше \(10^{-4}\tau\). Для устойчивого направления частоты \(e_i\) уже образуют единичный базис, так что запас в границе (1) сохраняется.

\end{document}
